\documentclass[11pt,a4paper]{article}
\usepackage{practicalreports}
\usepackage[margin=2.5cm]{geometry}
\usepackage{siunitx}
\usepackage{booktabs}

\setpagefancy
\header{Physics Lab --- Experiment 3}{}{\today}
\goodfoot

\begin{document}

\newtitlepage{Measuring $g$ with a Simple Pendulum}
{J. Okafor \quad Physics 121, Section B}
{\today}

\section{Objective}

To determine the local acceleration due to gravity, $g$, by measuring
the period of oscillation of a simple pendulum at several lengths and
fitting the small-angle relation
\[
T = 2\pi\sqrt{\frac{L}{g}}.
\]

\section{Method}

A steel bob (mass \SI{50}{\gram}) was suspended from a fixed clamp by
a length of inextensible string. For each of six string lengths $L$,
the bob was displaced through a small angle ($<10^\circ$, so the
small-angle approximation holds) and released from rest. The time for
20 complete oscillations was measured with a stopwatch and divided by
20 to give the period $T$, reducing the relative timing error from a
single human reaction time by a factor of 20. Each length was
measured three times and averaged.

\shift{0.5}{\textit{Apparatus note: string length $L$ was measured
from the point of suspension to the centre of the bob, not to the top
of the bob, since the pendulum's effective length is measured to its
centre of mass.}}

\section{Results}

\newtable{Period vs.\ pendulum length}
{ccc}
{$L$ (\si{\meter}) & $T$ (\si{\second}, mean of 3) & $T^2$ (\si{\second\squared})}
{
0.20 & 0.90 & 0.81 \\
0.40 & 1.27 & 1.61 \\
0.60 & 1.55 & 2.40 \\
0.80 & 1.80 & 3.24 \\
1.00 & 2.01 & 4.04 \\
1.20 & 2.20 & 4.84
}
{pendulum-data}

\reftable{pendulum-data} shows a clear linear relationship between
$L$ and $T^2$, consistent with the small-angle model. A linear
least-squares fit of $T^2$ against $L$ gives a slope of
$m = \SI{4.037}{\second\squared\per\meter}$. Since
$T^2 = \frac{4\pi^2}{g}L$, the slope is $4\pi^2/g$, so

\boxanswer{g = \frac{4\pi^2}{m} = \frac{4\pi^2}{4.037} = \SI{9.78}{\meter\per\second\squared}}

\section{Uncertainty analysis}

The dominant source of uncertainty is the stopwatch reaction time,
estimated at $\pm\SI{0.2}{\second}$ per start/stop. Dividing by 20
oscillations gives a period uncertainty of $\pm\SI{0.01}{\second}$ per
measurement. Propagating this through the linear fit gives a slope
uncertainty of $\pm\SI{0.06}{\second\squared\per\meter}$, and hence

\[
\delta g = g \cdot \frac{\delta m}{m} = 9.78 \times \frac{0.06}{4.037}
\approx \SI{0.15}{\meter\per\second\squared}.
\]

The length measurement (a metre rule, $\pm\SI{1}{\milli\meter}$)
contributes negligibly by comparison at the lengths used here.

\section{Discussion}

The measured value $g = \SI{9.78 \pm 0.15}{\meter\per\second\squared}$
agrees with the standard value $\SI{9.81}{\meter\per\second\squared}$
to within one standard uncertainty. The small systematic
undershoot is consistent with air resistance and small residual
amplitude effects, both of which slightly lengthen the measured
period relative to the ideal small-angle prediction and would bias
the fitted $g$ downward, exactly the direction of the observed
discrepancy.

\section{Conclusion}

The simple-pendulum method gives a result for $g$ consistent with
the accepted value, and the linear $T^2$-vs-$L$ relationship
predicted by the small-angle approximation is well supported by the
data across the full range of lengths tested.

\end{document}
